% ============================================================================
%  04-05: TRANSACCIONES, TRIGGERS E INDEXACIÓN
% ============================================================================

\section{Transacciones, triggers e indexaci\'on}

\subsection{Procedimiento transaccional \texttt{sp\_confirmar\_venta}}

El procedimiento almacenado \texttt{sp\_confirmar\_venta} encapsula todo el ciclo de confirmaci\'on de un pedido dentro de una transacci\'on at\'omica con bloqueo pesimista (\texttt{SELECT FOR UPDATE}).
Esto garantiza que dos sesiones concurrentes no puedan vender el mismo \'ultimo stock disponible.

\begin{lstlisting}[language=SQL, caption={Procedimiento transaccional con bloqueo pesimista}, label=lst:sp-venta]
CREATE OR REPLACE PROCEDURE ventas.sp_confirmar_venta(
  p_user_id      TEXT,
  p_items        JSONB,   -- [{componenteId, cantidad, precio}]
  p_metodo_pago  SMALLINT
)
LANGUAGE plpgsql
AS $$
DECLARE
  v_pedido_id  INTEGER;
  v_total      NUMERIC(12,2) := 0;
  v_item       JSONB;
  v_comp_id    INTEGER;
  v_cantidad   INTEGER;
  v_precio     NUMERIC(10,2);
  v_stock_act  INTEGER;
BEGIN
  -- 1. Bloqueo pesimista de filas a modificar
  FOR v_item IN SELECT * FROM jsonb_array_elements(p_items) LOOP
    v_comp_id  := (v_item->>'componenteId')::INTEGER;
    v_cantidad := (v_item->>'cantidad')::INTEGER;

    SELECT "Stock" INTO v_stock_act
    FROM catalogo."Componentes"
    WHERE "Id" = v_comp_id
    FOR UPDATE;  -- bloqueo exclusivo de fila

    IF v_stock_act < v_cantidad THEN
      RAISE EXCEPTION
        'Stock insuficiente para componente %. Disponible: %, Requerido: %',
        v_comp_id, v_stock_act, v_cantidad;
    END IF;
  END LOOP;

  -- 2. Crear cabecera del pedido
  INSERT INTO ventas."Pedidos"
    ("UserId", "Total", "MetodoPago", "Estado")
  VALUES
    (p_user_id, 0, p_metodo_pago, 0)
  RETURNING "Id" INTO v_pedido_id;

  -- 3. Insertar items y decrementar stock
  FOR v_item IN SELECT * FROM jsonb_array_elements(p_items) LOOP
    v_comp_id  := (v_item->>'componenteId')::INTEGER;
    v_cantidad := (v_item->>'cantidad')::INTEGER;
    v_precio   := (v_item->>'precio')::NUMERIC;

    INSERT INTO ventas."PedidoItems"
      ("PedidoId","ComponenteId","Cantidad","PrecioUnitario")
    VALUES
      (v_pedido_id, v_comp_id, v_cantidad, v_precio);

    UPDATE catalogo."Componentes"
       SET "Stock" = "Stock" - v_cantidad
     WHERE "Id"   = v_comp_id;

    v_total := v_total + (v_cantidad * v_precio);
  END LOOP;

  -- 4. Actualizar total del pedido
  UPDATE ventas."Pedidos"
     SET "Total" = v_total,
         "Estado" = 1  -- Confirmado
   WHERE "Id" = v_pedido_id;

  COMMIT;
EXCEPTION
  WHEN OTHERS THEN
    ROLLBACK;
    RAISE;
END;
$$;
\end{lstlisting}

\subsection{Flujo de la transacci\'on}

\figuraAPA{fig:flujo-tx}{Flujo de la transacci\'on \texttt{sp\_confirmar\_venta}}{%
\begin{tikzpicture}[
  nodo/.style={draw, rectangle, rounded corners=4pt, fill=white, minimum width=3.5cm, minimum height=0.7cm, font=\small, align=center},
  dec/.style={draw, diamond, fill=yellow!20, minimum width=3.0cm, minimum height=0.7cm, font=\small, aspect=2.5, align=center},
  ok/.style={-Stealth, thick, green!50!black},
  err/.style={-Stealth, thick, red},
  norm/.style={-Stealth, thick},
  node distance=0.9cm
]
\node[nodo, fill=blue!10] (ini) {BEGIN TRANSACTION};
\node[nodo, below=of ini] (lock) {SELECT FOR UPDATE\\(bloqueo pesimista)};
\node[dec, below=of lock] (chk) {Stock suficiente?};
\node[nodo, right=2.5cm of chk, fill=red!10] (rb) {ROLLBACK\\RAISE EXCEPTION};
\node[nodo, below=of chk] (ins) {INSERT Pedido\\INSERT PedidoItems};
\node[nodo, below=of ins] (upd) {UPDATE Stock\\UPDATE Total};
\node[nodo, below=of upd, fill=green!10] (commit) {COMMIT};

\draw[norm] (ini)   -- (lock);
\draw[norm] (lock)  -- (chk);
\draw[ok]   (chk)   -- node[left,font=\tiny]{Sí} (ins);
\draw[err]  (chk)   -- node[above,font=\tiny]{No} (rb);
\draw[norm] (ins)   -- (upd);
\draw[norm] (upd)   -- (commit);
\end{tikzpicture}%
}{Elaboraci\'on propia}

\subsection{Triggers de auditor\'ia}

Los triggers de auditor\'ia registran cada inserci\'on, actualizaci\'on y eliminaci\'on en las tablas cr\'iticas, proporcionando un historial inmutable de cambios.

\begin{lstlisting}[language=SQL, caption={Trigger de auditoría para Componentes}, label=lst:trigger]
-- Tabla de auditoría
CREATE TABLE auditoria."AuditoriaComponentes" (
  "Id"         SERIAL PRIMARY KEY,
  "Operacion"  CHAR(1)     NOT NULL, -- I=Insert, U=Update, D=Delete
  "ComponenteId" INTEGER,
  "StockAntes" INTEGER,
  "StockDespues" INTEGER,
  "PrecioAntes" NUMERIC(10,2),
  "PrecioDespues" NUMERIC(10,2),
  "FechaHora"  TIMESTAMPTZ NOT NULL DEFAULT NOW(),
  "Usuario"    TEXT        NOT NULL DEFAULT current_user
);

-- Función del trigger
CREATE OR REPLACE FUNCTION auditoria.fn_audit_componentes()
RETURNS TRIGGER LANGUAGE plpgsql AS $$
BEGIN
  IF TG_OP = 'INSERT' THEN
    INSERT INTO auditoria."AuditoriaComponentes"
      ("Operacion","ComponenteId","StockDespues","PrecioDespues")
    VALUES ('I', NEW."Id", NEW."Stock", NEW."Precio");
  ELSIF TG_OP = 'UPDATE' THEN
    INSERT INTO auditoria."AuditoriaComponentes"
      ("Operacion","ComponenteId",
       "StockAntes","StockDespues",
       "PrecioAntes","PrecioDespues")
    VALUES ('U', NEW."Id",
      OLD."Stock", NEW."Stock",
      OLD."Precio", NEW."Precio");
  ELSIF TG_OP = 'DELETE' THEN
    INSERT INTO auditoria."AuditoriaComponentes"
      ("Operacion","ComponenteId","StockAntes","PrecioAntes")
    VALUES ('D', OLD."Id", OLD."Stock", OLD."Precio");
  END IF;
  RETURN NEW;
END;
$$;

-- Binding del trigger
CREATE TRIGGER tg_audit_componentes
AFTER INSERT OR UPDATE OR DELETE
ON catalogo."Componentes"
FOR EACH ROW EXECUTE FUNCTION auditoria.fn_audit_componentes();
\end{lstlisting}

\subsection{Estrategia de indexaci\'on}

\begin{center}
\begin{tablaAPA}
\begin{tabular}{L{3.5cm} L{3.5cm} C{1.5cm} L{4.5cm}}
\toprule
\textbf{\'Indice} & \textbf{Tabla/Columna} & \textbf{Tipo} & \textbf{Justificaci\'on} \\
\midrule
idx\_comp\_categoria & Componentes.CategoriaId & B-Tree & Filtros de categor\'ia en la tienda \\
idx\_comp\_marca     & Componentes.MarcaId     & B-Tree & Filtros de marca en la tienda \\
idx\_comp\_precio    & Componentes.Precio      & B-Tree & Ordenamiento y rangos de precio \\
idx\_pedidos\_user   & Pedidos.UserId          & B-Tree & Historial de pedidos por cliente \\
idx\_pedidos\_estado & Pedidos.Estado          & B-Tree & Filtro por estado en panel admin \\
idx\_pedidos\_fecha  & Pedidos.FechaPedido     & B-Tree & Consultas de per\'iodo (reportes) \\
idx\_carritos\_json  & CarritosAbandonados.ItemsJson & GIN & B\'usqueda dentro del JSONB \\
\bottomrule
\end{tabular}
\end{tablaAPA}
\captionof{table}{\'Indices definidos en el esquema DESPU\'ES}
\label{tab:indices}
\end{center}

\begin{lstlisting}[language=SQL, caption={Creación de índices B-Tree y GIN}, label=lst:indices]
CREATE INDEX idx_comp_categoria
  ON catalogo."Componentes"("CategoriaId");
CREATE INDEX idx_comp_marca
  ON catalogo."Componentes"("MarcaId");
CREATE INDEX idx_comp_precio
  ON catalogo."Componentes"("Precio");
CREATE INDEX idx_pedidos_user
  ON ventas."Pedidos"("UserId");
CREATE INDEX idx_pedidos_estado
  ON ventas."Pedidos"("Estado");
CREATE INDEX idx_pedidos_fecha
  ON ventas."Pedidos"("FechaPedido" DESC);
CREATE INDEX idx_carritos_json
  ON ventas."CarritosAbandonados"
  USING GIN ("ItemsJson");
\end{lstlisting}

\subsection{An\'alisis de plan de ejecuci\'on}

La siguiente tabla compara el tiempo de ejecuci\'on estimado para la consulta de productos por categor\'ia antes y despu\'es de crear el \'indice \texttt{idx\_comp\_categoria}.

\begin{center}
\begin{tablaAPA}
\begin{tabular}{L{4.0cm} C{3.5cm} C{3.5cm} L{2.0cm}}
\toprule
\textbf{Escenario} & \textbf{Plan utilizado} & \textbf{Costo estimado} & \textbf{Mejora} \\
\midrule
Sin \'indice (Seq Scan) & Sequential Scan & 85.42 ms & --- \\
Con \'indice (Index Scan) & Index Scan Backward & 0.08 ms & 99.9\% \\
\bottomrule
\end{tabular}
\end{tablaAPA}
\captionof{table}{Comparativa \texttt{EXPLAIN ANALYZE} antes y despu\'es de indexaci\'on}
\label{tab:explain}
\end{center}

La reducci\'on del 99.9\% en el tiempo de ejecuci\'on se debe a que PostgreSQL puede saltar directamente a las filas de la categor\'ia buscada usando el \'indice B-Tree, evitando el escaneo secuencial de toda la tabla \texttt{Componentes}.
